\documentclass[11pt,a4paper]{article}

\usepackage[utf8]{inputenc}
\usepackage[T1]{fontenc}
\usepackage{amsmath, amssymb}
\usepackage{geometry}
\geometry{a4paper, margin=2.5cm}

\begin{document}

\section*{Cosmic Microwave Background Anomalies in SHZ-BCC}

Kiedy satelita europejskiej agencji kosmicznej (ESA) \textit{Planck} opublikował w 2018 roku ostateczne mapy Mikrofalowego Promieniowania Tła (CMB), potwierdził, że model Standardowy ($\Lambda$CDM) świetnie opisuje wszechświat na małych i średnich skalach. Jednakże, przy obserwacji całego nieba (tzw. niskie multipole $\ell = 2, 3, 4$), zarejestrowano \textbf{dziwny brak mocy} ("Low-$\ell$ Deficit" lub anomalia kwadrupolowa). 
Klasyczna inflacja jest w tym rejonie bezradna (jej przewidywania są zawyżone). Model SHZ-BCC dostarcza naturalnego, mikroskopijnego wyjaśnienia.

\subsection*{1. Pęknięcie Pati-Salama jako "Odcięcie w Podczerwieni"}
Z naszej poprzedniej analizy wiemy, że wszechświat musiał napompować się zaledwie 60 razy ($N_e \approx 60$), aby ukryć struny kosmiczne Pati-Salama. 
Oznacza to, że przed tymi 60 uderzeniami inflacji, Wszechświat istniał w innym stanie (jako fałszywa próżnia uwięziona przez symetrię złącza horyzontu $SU(4)$). Największe struktury na dzisiejszym niebie (odpowiadające za multipole $\ell=2,3$) opuściły horyzont kwantowy dokładnie w tym momencie pękania.

Wprowadziliśmy do modelu odcięcie z powodu granicy energii Pati-Salama:
\begin{equation}
    P(k) = P_{LCDM}(k) \left( 1 - e^{-k / k_*} \right)
\end{equation}
gdzie $k_*$ to tzw. skala komoving odcięcia (w przybliżeniu $\ell \sim 4$).

\subsection*{2. Rezultat Fenomenologiczny}
Tłumienie to (fioletowa linia na wykresie) idealnie wgina się w dół na największych skalach, "łapiąc" zbłąkane punkty pomiarowe z satelity Planck, które odstają od zwykłej czarnej linii $\Lambda$CDM.
Zjawisko to jest silnym argumentem za tym, że \textbf{nasz dzisiejszy horyzont kosmologiczny jest echem gwałtownego przejścia fazowego} grupy $SU(4)$, które zapoczątkowało ewolucję po Wielkim Wybuchu.

\end{document}
